% Note this requires the (standard) exam package
\documentclass[a4paper,12pt,addpoints]{exam} 

%This is the package which you need to include. The ``onlynumbers'' option make part numbering numerical only (e.g. 2.1.5)
%\usepackage[onlynumbers]{buexam}
% By default, part numbering is number.letter.roman, e.g. 2.a.v
\usepackage{buexam}



% The Unit title of the exam (append RESIT to the end if necessary)
\setunittitle{Zombies in Popular Media and Mathematics}
% A shortened / abbreviated name of the unit
\setunittitleshort{ZPMM}
% The unit reference (might be the same as short name)
\setunitrefs{ZPMM}
% The level of the examination (4,5,6 etc)
\setunitlevel{L4}
% The exam reference (get this from the school
\setexamref{ZZZ1234Z}
% The date of the exam
\setexamdate{10/06/2013}
% Unit leader of the course
\setunitleadername{Anne Onymous}
% Unit leader extension
\setunitleaderext{60000}
% Assistant name
\setassistantname{John Doe}
% Assitant extension
\setassistantext{60000}


% There are the specific instructions for candidates
\setinstructions{
See Front Page Instruction Samples document \href{https://sasportal.bournemouth.ac.uk/exams/Exam\%20Question\%20Paper\%20Documents/Example\%20INSTRUCTIONS\%20(exam\%20papers).docx}{[here]}.

This could include information about:
\begin{itemize}[leftmargin=0pt]
\item[] If question paper needs to be handed in
\item[] Questions to answer
\item[] Allocation of marks
\item[] Where / how to answer questions (e.g. answer books /  MCQ answer sheet etc)
\item[] Supplied materials
\item[] Permitted materials / equipment
\item[] Note, this list should be created without bullet points (hence the \textbackslash item[] and the 0pt leftmargin)
\end{itemize}
}

% This option should be commented out for the final paper (don't forget!)
%\printanswers

\begin{document}

% Note that this command will create ``Section B\\Answer All Questions'' - you have to use this, as it is in an itemized environment.
\sectioninstructions{Answer ALL Questions}

% Questions can consist of multiple parts. Note that the sum of the parts must make the total (in this case, [15]).
\titledquestion{Zombies in films}
Consider Zombies in films. 
\begin{parts}
\part[5] Name some films with Zombies \droppoints
\begin{solution}
\begin{itemize}
\item Pride and Prejudice... with Zombies!
\item Abraham Lincoln... Zombie Hunter!
\item etc.
\end{itemize}
\end{solution}

% Maths example
\part[10] State the average number of characters eaten by a single Zombie in "`28 Days Later"'. Use the following formula: \droppoints
\begin{align}
\bar{x} &= \frac{1}{n}\cdot \sum_{i=1}^{n} x_i \nonumber
\end{align}
\begin{solution}
At least 3. Zombies are awesome.
\end{solution}
\end{parts}

% A newpage is required between sections
\newpage

% Note that this command will create ``Section B\\Answer Only ODD numbered questions'' - you have to use this, as it is in an itemized environment.
\sectioninstructions{Answer Only ODD numbered questions}

\titledquestion{Zombie Programming}
Consider the following code segment:
\begin{CPP}
#include <iostream>
using namespace std;

int main ()
{
  cout << "Hello World!";
  return 0;
}
\end{CPP}
\begin{parts}
%Example of code inclusion
\part In your own words, explain:
\begin{subparts}
\subpart[5] Was this code produced by a Zombie? If so, why?\droppoints
\begin{solution}
Clearly, the lameness indicates that it is Zombie generated code (ZGC).
\end{solution}
\subpart[5] How would you improve on the above code? \droppoints 
\begin{solution}
Don't use the std namespace! Better to give your fingers exercise and type it every time.
\end{solution}
\end{subparts}

\part[10] Write a page essay on any topic you like. Appropriate use of references required.\droppoints
\begin{solution}
 Blah Blah Blah
\end{solution}

\end{parts}

\newpage

\titledquestion{Zombie Algorithms}
\begin{parts}

% Example of an algorithm
 \part[10] Consider the following algorithm:

\begin{algorithmic}
\REQUIRE $n \geq 0$
\ENSURE $y = x^n$
\STATE $y \leftarrow 1$
\STATE $X \leftarrow x$
\STATE $N \leftarrow n$
\WHILE{$N \neq 0$}
\IF{$N$ is even}
\STATE $X \leftarrow X \times X$
\STATE $N \leftarrow N / 2$
\ELSE[$N$ is odd]
\STATE $y \leftarrow y \times X$
\STATE $N \leftarrow N - 1$
\ENDIF
\ENDWHILE
\end{algorithmic}
Describe in your own words the thought processes required for a Zombie to generate this algorithm. How can the promise of fresh brains serve as a motivators? \droppoints
\begin{solution}
Blah blah blah.
\end{solution}

% Example of image inclusion
\part[5] Admire our institution's mighty logo:
\begin{center}
\includegraphics[width=0.1\textwidth]{bulogo.pdf}
\end{center}
Draw this in a manner similar to a Zombie. Coloured pens are provided.\droppoints
\begin{solution}
Blah blah blah.
\end{solution}
\end{parts}

% The postamble is needed as part of the uni style. (this cannot be automatic as AtEndOfDocument adds a \newpage)
\afterlastquestion

\end{document}


